\documentclass[UTF8]{ctexart}
\usepackage{xeCJK}
\usepackage{geometry}
\geometry{left=2cm,right=2cm,top=2.5cm,bottom=2.5cm}
\setlength{\headheight}{12.7pt}

\usepackage{titlesec}
\titleformat{\section}
  {\normalfont\large\bfseries}
  {\thesection}
  {1em}
  {}
\titlespacing*{\section}
  {0pt}
  {3.5ex plus 1ex minus .2ex}
  {2.3ex plus .2ex}

\usepackage{amsmath}
\usepackage{amssymb}
\usepackage{amsthm}
\usepackage{enumitem}
\setlist[enumerate]{itemsep=0pt,topsep=0pt,parsep=0pt,leftmargin=0pt,itemindent=2.5em}
\setlist[itemize]{itemsep=0pt,topsep=0pt,parsep=0pt,leftmargin=0pt,itemindent=2.5em}
\usepackage{fancyhdr}
\pagestyle{fancy}
\fancyhf{}
\usepackage{graphicx}
\usepackage{hyperref}
\usepackage{indentfirst}
\usepackage{lmodern}

\begin{document}
\setlength{\abovedisplayskip}{1pt}
\setlength{\belowdisplayskip}{1pt}

\section{一阶语言中的替换}

\hypertarget{sec:定义1.1}{}
\noindent\textbf{定义1.1 (替换)}

给定一阶语言$\mathcal{L}$,
称映射$\theta:\mathcal{V}\cup\mathcal{C}\to T^{\mathcal{L}}$为一个
$\mathcal{L}$-\textbf{替换}. 把嵌入映射称作恒等替换, 记为$i$.

\vspace{10pt}

此时再给定变元$x$和项$t$, 定义替换$\theta$的单点修改
$\displaystyle\theta\frac{t}{x}$如下:
\[\forall y\in\mathcal{V}\cup\mathcal{C}:\quad\theta\frac{t}{x}(y)=
\left\{\begin{aligned}
    &t,\quad&y=x,\\[-3pt]
    &\theta(y),\quad&y\neq x,
\end{aligned}\right.\]
即把$\theta$对$x$的替换修改成$t$, 得到的新的替换.
把$\displaystyle\theta\frac{x}{x}$简记为$\theta/x$. 显然总有$i/x=i$.

\vspace{10pt}

\hypertarget{sec:定义1.2}{}
\noindent\textbf{定义1.2 (项的替换)}

给定一阶语言$\mathcal{L}$和$\mathcal{L}$-替换$\theta$.
在$T^{\mathcal{L}}$中定义项的映射$\widehat{\theta}$:
\begin{enumerate}
    \item 对任意的原子项$t$, $\widehat{\theta}(t)=\theta(t)$,
    \item 对任意的$n$元函数符号$f$和项$t_1,\cdots,t_n$,
    $\widehat{\theta}(ft_1\cdots t_n)=f\widehat{\theta}(t_1)\cdots
    \widehat{\theta}(t_n)$.
\end{enumerate}
在不产生歧义的情况下通常把$\widehat{\theta}(t)$直接写成$\theta(t)$, 称为$\theta$对项$t$的替换.

\vspace{10pt}

\hypertarget{sec:定义1.3}{}
\noindent\textbf{定义1.3 (公式的替换)}

给定一阶语言$\mathcal{L}$.
在$F^{\mathcal{L}}$中(同时对所有替换$\theta$)定义公式的映射$\overline{\theta}$:
\begin{enumerate}
    \item 对任意的项$t_1,t_2$,
    $\overline{\theta}(t_1\equiv t_2)=\widehat{\theta}(t_1)\equiv
    \widehat{\theta}(t_2)$,
    \item 对任意的$n$元关系符号$r$和项$t_1,\cdots,t_n$,
    $\overline{\theta}(rt_1\cdots t_n)=r \widehat{\theta}(t_1)\cdots
    \widehat{\theta}(t_n)$,
    \item 对任意的公式$\varphi$,
    $\overline{\theta}(\neg\varphi)=\neg \overline{\theta}(\varphi)$,
    \item 对任意的公式$\varphi,\psi$, $\overline{\theta}((\varphi\to\psi))
    =(\overline{\theta}(\varphi)\to \overline{\theta}(\psi))$,
    \item 对任意的变元$x$和公式$\varphi$,
    $\overline{\theta}(\exists x\varphi)=\exists x\,\overline{\theta/x}(\varphi)$.
\end{enumerate}
在不产生歧义的情况下通常把$\overline{\theta}(\varphi)$直接写成$\theta(\varphi)$,
称为$\theta$对公式$\varphi$的替换.

\vspace{10pt}

易证恒等替换$i$保持所有的项和公式不变.

\vspace{10pt}

\hypertarget{sec:定义1.4}{}
\noindent\textbf{定义1.4 (自由替换)}

给定一阶语言$\mathcal{L}$, (同时对所有替换$\theta$)定义公式的命题$R^\theta$:
\begin{enumerate}
    \item 对任意的原子公式$\varphi$, $R^\theta(\varphi)$为真,
    \vspace{3pt}
    \item 对任意的公式$\varphi$,
    $R^\theta(\neg\varphi)\overset{\mathrm{def}}{\leftrightarrow}R^\theta(\varphi)$,
    \vspace{3pt}
    \item 对任意的公式$\varphi,\psi$,
    $R^\theta((\varphi\to\psi))\overset{\mathrm{def}}{\leftrightarrow}
    R^\theta(\varphi)\land R^\theta(\psi)$,
    \vspace{3pt}
    \item 对任意的变元$x$和公式$\varphi$,
    $R^\theta(\exists x\varphi)\overset{\mathrm{def}}{\leftrightarrow}
    R^{\theta/x}(\varphi)
    \land\forall y\in\mathrm{at}(\exists x\varphi):x\notin\mathrm{var}\,\theta(y)$.
\end{enumerate}
把$R^\theta(\varphi)$记作$\theta[\varphi]$,
这时称替换$\theta$对公式$\varphi$是\textbf{自由的}.

\vspace{10pt}

易证恒等替换$i$对所有的公式都是自由的.





\newpage

\section{替换的重合引理}

\hypertarget{sec:定理2.1}{}
\noindent\textbf{定理2.1}

给定一阶语言$\mathcal{L}$. 对任意的项$t$和替换$\theta,\eta$, 如果
\[
    \forall x\in\mathrm{at}(t):\theta(x)=\eta(x),
\]
记作$\theta\overset{\displaystyle t}{=}\eta$, 那么$\theta(t)=\eta(t)$.

\noindent{\kaishu 证明.}

\begin{enumerate}
    \item 显然原子项$t$的条件与结论相同.
    \item 对任意的$n$元函数符号$f$, 假设$t_1,\cdots,t_n$满足命题并且
    $t=ft_1\cdots t_n$满足$\theta\overset{\displaystyle t}{=}\eta$,
    则对每个$t_i$有
    \[
        \mathrm{at}(t_i)\subset\mathrm{at}(t)\quad\implies\quad
        \theta\overset{\displaystyle t_i}{=}\eta\quad\implies\quad
        \theta(t_i)=\eta(t_i),
    \]
    从而$\theta(t)=f\theta(t_1)\cdots\theta(t_n)=f\eta(t_1)\cdots\eta(t_n)=\eta(t)$.
    \hfill $\qedsymbol$
\end{enumerate}

\vspace{10pt}

\hypertarget{sec:定理2.2}{}
\noindent\textbf{定理2.2}

给定一阶语言$\mathcal{L}$. 对任意的公式$\varphi$和替换$\theta,\eta$, 如果
\[
    \forall x\in\mathrm{at}(\varphi):\theta(x)=\eta(x),
\]
记作$\theta\overset{\displaystyle\varphi}{=}\eta$,
那么$\theta(\varphi)=\eta(\varphi)$.

\noindent{\kaishu 证明.}

\begin{enumerate}
    \item 对任意的项$t_1,t_2$, 假设$\varphi=t_1\equiv t_2$满足
    $\theta\overset{\displaystyle\varphi}{=}\eta$, 则对每个$t_i$有
    \[
        \mathrm{at}(t_i)\subset\mathrm{at}(\varphi)\quad\implies\quad
        \theta\overset{\displaystyle t_i}{=}\eta\quad\implies\quad
        \theta(t_i)=\eta(t_i),
    \]
    从而$\theta(\varphi)=\theta(t_1)\equiv\theta(t_2)=
    \eta(t_1)\equiv\eta(t_2)=\eta(\varphi)$.

    \item 对任意的$n$元关系符号$r$和项$t_1,\cdots,t_n$, 假设$\varphi=rt_1\cdots t_n$
    满足$\theta\overset{\displaystyle\varphi}{=}\eta$,
    同理每个$\theta(t_i)=\eta(t_i)$, 从而
    \[
        \theta(\varphi)=r\theta(t_1)\cdots\theta(t_n)=
        r\eta(t_1)\cdots\eta(t_n)=\eta(\varphi).
    \]

    \item 假设$\varphi_1$满足命题并且$\varphi=\neg\varphi_1$满足
    $\theta\overset{\displaystyle\varphi}{=}\eta$.
    因为$\mathrm{at}(\varphi)=\mathrm{at}(\varphi_1)$, 所以
    $\theta\overset{\displaystyle\varphi_1}{=}\eta$, 从而
    \[
        \theta(\varphi)=\neg\theta(\varphi_1)=\neg\eta(\varphi_1)=\eta(\varphi).
    \]

    \item 假设$\varphi_1,\varphi_2$满足命题并且$\varphi=(\varphi_1\to\varphi_2)$满足
    $\theta\overset{\displaystyle\varphi}{=}\eta$.

    \hspace{2.17em}
    因为$\mathrm{at}(\varphi)=\mathrm{at}(\varphi_1)\cup\mathrm{at}(\varphi_2)$,
    所以$\theta\overset{\displaystyle\varphi_1}{=}\eta$并且
    $\theta\overset{\displaystyle\varphi_2}{=}\eta$, 从而
    \[
        \theta(\varphi)=(\theta(\varphi_1)\to\theta(\varphi_2))=
        (\eta(\varphi_1)\to\eta(\varphi_2))=\eta(\varphi).
    \]

    \item 对任意的变元$x$, 假设$\varphi_1$满足命题并且$\varphi=\exists x\varphi_1$
    满足$\theta\overset{\displaystyle\varphi}{=}\eta$.
    下证$\theta/x\overset{\displaystyle\varphi_1}{=}\eta/x$,
    任取$y\in\mathrm{at}(\varphi_1)$.

    \hspace{2.17em}
    如果$y=x$, 那么$\theta/x(y)=x=\eta/x(y)$; 如果$y\neq x$,
    那么$y\in\mathrm{at}(\varphi)$, 依假设有
    \[
        \theta/x(y)=\theta(y)=\eta(y)=\eta/x(y),
    \]
    得证. 依假设有$\theta/x(\varphi_1)=\eta/x(\varphi_1)$, 从而
    \begin{equation*}
        \theta(\varphi)=\exists x(\theta/x)(\varphi_1)=
        \exists x(\eta/x)(\varphi_1)=\eta(\varphi).
    \tag*{\qedsymbol}\end{equation*}
\end{enumerate}





\newpage

\section{替换后的变元和常元}

\hypertarget{sec:定理3.1}{}
\noindent\textbf{定理3.1}

给定一阶语言$\mathcal{L}$. 对任意的项$t$和替换$\theta$, 有$\displaystyle
\mathrm{at}\,\theta(t)=\bigcup_{x\in\mathrm{at}(t)}\mathrm{at}\,\theta(x)$.

\noindent{\kaishu 证明.}

\begin{enumerate}
    \item 对任意的原子项$t$, 由$\mathrm{at}(t)=\{t\}$即得.
    \item 对任意的$n$元函数符号$f$, 假设$t_1,\cdots,t_n$满足命题, 那么
    对$t=ft_1\cdots t_n$有
    \vspace{3pt}
    \begin{equation*}
        \mathrm{at}\,\theta(t)=
        \mathrm{at}\,\theta(t_1)\cup\cdots\cup\mathrm{at}\,\theta(t_n)=
        \left(\bigcup_{x\in\mathrm{at}(t_1)}\mathrm{at}\,\theta(x)\right)\cup\cdots
        \cup\left(\bigcup_{x\in\mathrm{at}(t_n)}\mathrm{at}\,\theta(x)\right)=
        \bigcup_{x\in\mathrm{at}(t)}\mathrm{at}\,\theta(x).
    \tag*{\qedsymbol}\end{equation*}
\end{enumerate}

\vspace{10pt}

\hypertarget{sec:定理3.2}{}
\noindent\textbf{定理3.2}

给定一阶语言$\mathcal{L}$. 对任意的公式$\varphi$和替换$\theta$, 如果
$\theta[\varphi]$, 那么$\displaystyle\mathrm{at}\,\theta(\varphi)=
\bigcup_{x\in\mathrm{at}(\varphi)}\mathrm{at}\,\theta(x)$.

\noindent{\kaishu 证明.}

\begin{enumerate}
    \item 对任意的项$t_1,t_2$, 对$\varphi=t_1\equiv t_2$有
    \[
        \mathrm{at}\,\theta(\varphi)=
        \mathrm{at}\,\theta(t_1)\cup\mathrm{at}\,\theta(t_2)=
        \left(\bigcup_{x\in\mathrm{at}(t_1)}\mathrm{at}\,\theta(x)\right)\cup
        \left(\bigcup_{x\in\mathrm{at}(t_2)}\mathrm{at}\,\theta(x)\right)=
        \bigcup_{x\in\mathrm{at}(\varphi)}\mathrm{at}\,\theta(x).
    \]

    \vspace{3pt}
    \item 对任意的$n$元关系符号$r$和项$t_1,\cdots,t_n$, 对$\varphi=rt_1\cdots t_n$有
    \vspace{3pt}
    \[
        \mathrm{at}\,\theta(\varphi)=
        \mathrm{at}\,\theta(t_1)\cup\cdots\cup\mathrm{at}\,\theta(t_n)=
        \left(\bigcup_{x\in\mathrm{at}(t_1)}\mathrm{at}\,\theta(x)\right)\cup\cdots
        \cup\left(\bigcup_{x\in\mathrm{at}(t_n)}\mathrm{at}\,\theta(x)\right)=
        \bigcup_{x\in\mathrm{at}(\varphi)}\mathrm{at}\,\theta(x).
    \]
    
    \vspace{3pt}
    \item 假设$\varphi_1$满足命题并且$\varphi=\neg\varphi_1$满足$\theta[\varphi]$,
    此时也有$\theta[\varphi_1]$, 从而
    \vspace{3pt}
    \[
        \mathrm{at}\,\theta(\varphi)=\mathrm{at}\,\theta(\varphi_1)=
        \bigcup_{x\in\mathrm{at}(\varphi_1)}\mathrm{at}\,\theta(x)=
        \bigcup_{x\in\mathrm{at}(\varphi)}\mathrm{at}\,\theta(x).
    \]

    \vspace{3pt}
    \item 假设$\varphi_1,\varphi_2$满足命题并且$\varphi=(\varphi_1\to\varphi_2)$
    满足$\theta[\varphi]$, 此时也有$\theta[\varphi_1],\theta[\varphi_2]$, 从而
    \[
        \mathrm{at}\,\theta(\varphi)=
        \mathrm{at}\,\theta(\varphi_1)\cup\mathrm{at}\,\theta(\varphi_2)=
        \left(\bigcup_{x\in\mathrm{at}(\varphi_1)}\mathrm{at}\,\theta(x)\right)\cup
        \left(\bigcup_{x\in\mathrm{at}(\varphi_2)}\mathrm{at}\,\theta(x)\right)=
        \bigcup_{x\in\mathrm{at}(\varphi)}\mathrm{at}\,\theta(x).
    \]

    \vspace{3pt}
    \item 对任意的变元$x$, 假设$\varphi_1$满足命题并且$\varphi=\exists x\varphi_1$
    满足$\theta[\varphi]$, 此时也有$\theta/x[\varphi_1]$, 从而
    \[
        \mathrm{at}\,\theta(\varphi)=
        \mathrm{at}\,(\theta/x)(\varphi_1)\setminus\{x\}=
        \left(\bigcup_{y\in\mathrm{at}(\varphi_1)}\mathrm{at}\,(\theta/x)(y)\right)
        \setminus\{x\}.
    \]

    \hspace{2.17em}
    如果$x\in\mathrm{at}(\varphi_1)$,
    此时$\mathrm{at}(\varphi_1)=\mathrm{at}(\varphi)\cup\{x\}$, 那么
    \vspace{3pt}
    \[
        \mathrm{at}\,\theta(\varphi)=
        \left(\bigcup_{y\in\mathrm{at}(\varphi)}\mathrm{at}\,(\theta/x)(y)\right)
        \cup\mathrm{at}\,(\theta/x)(x)\setminus\{x\}=
        \bigcup_{y\in\mathrm{at}(\varphi)}\mathrm{at}\,\theta(y).
    \]

    \hspace{2.17em}
    如果$x\notin\mathrm{at}(\varphi_1)$,
    此时$\mathrm{at}(\varphi_1)=\mathrm{at}(\varphi)$, 那么
    \[
        \mathrm{at}\,\theta(\varphi)=
        \left(\bigcup_{y\in\mathrm{at}(\varphi)}\mathrm{at}\,\theta(y)\right)
        \setminus\{x\},\\[3pt]
    \]
    因为$\theta[\varphi]$要求对任意的$y\in\mathrm{at}(\varphi)$有
    $x\notin\mathrm{var}\,\theta(y)$, 所以
    \begin{equation*}
        x\notin\bigcup_{y\in\mathrm{at}(\varphi)}\mathrm{at}\,\theta(y)\quad
        \implies\quad\mathrm{at}\,\theta(\varphi)=
        \bigcup_{y\in\mathrm{at}(\varphi)}\mathrm{at}\,\theta(y).
    \tag*{\qedsymbol}\end{equation*}
\end{enumerate}





\newpage

% 过了许久才发现, 这里还需要补充一些细节, 要说明某个地方的公式替换两次会回到自身.
% 还好从这个细节上发现了下面的还算有趣的结论, 算是有意外收获.

% 关于替换的引理太多了. 现在两个笔记加起来一共有 7 页.

\section{替换的复合}

为了避免混淆替换本身含义和其衍生的含义,
这里严格使用$\widehat{\theta},\overline{\theta}$这种符号.

\vspace{10pt}

\hypertarget{sec:定义4.1}{}
\noindent\textbf{定义4.1 (替换的复合)}

给定一阶语言$\mathcal{L}$. 对任意的替换$\theta,\eta$, 定义它们的\textbf{复合}
$\eta\circ\theta:y\mapsto \widehat{\eta}(\theta(y))$.

\vspace{10pt}

\hypertarget{sec:定理4.1}{}
\noindent\textbf{定理4.1}

给定一阶语言$\mathcal{L}$. 对任意的项$t$和替换$\theta,\eta$,
有$\widehat{\eta\circ\theta}(t)=\widehat{\eta}\widehat{\theta}(t)$.

\noindent{\kaishu 证明.}

\begin{enumerate}
    \item 对任意的原子项$t$, $\widehat{\eta\circ\theta}(t)=(\eta\circ\theta)(t)=
    \widehat{\eta}(\theta(t))=\widehat{\eta}\widehat{\theta}(t)$.
    \item 对任意的$n$元函数符号$f$, 假设$t_1,\cdots,t_n$满足命题,
    那么对$t=ft_1\cdots t_n$有
    \vspace{3pt}
    \begin{equation*}
        \widehat{\eta\circ\theta}(t)=
        f\widehat{\eta\circ\theta}(t_1)\cdots\widehat{\eta\circ\theta}(t_n)=
        f\widehat{\eta}\,\widehat{\theta}(t_1)\cdots
        \widehat{\eta}\,\widehat{\theta}(t_n)=
        \widehat{\eta}\left(f\widehat{\theta}(t_1)\cdots\widehat{\theta}(t_n)\right)
        =\widehat{\eta}\,\widehat{\theta}(t).
    \tag*{\qedsymbol}\end{equation*}
\end{enumerate}

\vspace{10pt}

\hypertarget{sec:定理4.2}{}
\noindent\textbf{定理4.2}

给定一阶语言$\mathcal{L}$. 对任意的公式$\varphi$和替换$\theta,\eta$,
如果$\theta[\varphi]$, 那么$\overline{\eta\circ\theta}(\varphi)=
\overline{\eta}\,\overline{\theta}(\varphi)$.

\noindent{\kaishu 证明.}

\begin{enumerate}
    \item 对任意的项$t_1,t_2$, 对$\varphi=t_1\equiv t_2$有
    \vspace{3pt}
    \[
        \overline{\eta\circ\theta}(\varphi)=
        \widehat{\eta\circ\theta}(t_1)\equiv\widehat{\eta\circ\theta}(t_2)=
        \widehat{\eta}\,\widehat{\theta}(t_1)\equiv
        \widehat{\eta}\,\widehat{\theta}(t_2)=
        \overline{\eta}\left(\widehat{\theta}(t_1)\equiv\widehat{\theta}(t_2)\right)
        =\overline{\eta}\,\overline{\theta}(\varphi).
    \]

    \vspace{3pt}
    \item 对任意的$n$元关系符号$r$和项$t_1,\cdots,t_n$, 对$\varphi=rt_1\cdots t_n$有
    \vspace{3pt}
    \[
        \overline{\eta\circ\theta}(\varphi)=
        r\widehat{\eta\circ\theta}(t_1)\cdots\widehat{\eta\circ\theta}(t_n)=
        r\widehat{\eta}\,\widehat{\theta}(t_1)\cdots
        \widehat{\eta}\,\widehat{\theta}(t_n)=
        \overline{\eta}\left(r\widehat{\theta}(t_1)\cdots\widehat{\theta}(t_n)\right)
        =\overline{\eta}\,\overline{\theta}(\varphi).
    \]
    
    \vspace{3pt}
    \item 假设$\varphi_1$满足命题并且$\varphi=\neg\varphi_1$满足$\theta[\varphi]$,
    此时也有$\theta[\varphi_1]$, 从而
    \vspace{3pt}
    \[
        \overline{\eta\circ\theta}(\varphi)=
        \neg\,\overline{\eta\circ\theta}(\varphi_1)=
        \neg\,\overline{\eta}\,\overline{\theta}(\varphi_1)=
        \overline{\eta}\left(\neg\,\overline{\theta}(\varphi_1)\right)=
        \overline{\eta}\,\overline{\theta}(\varphi).
    \]

    \vspace{3pt}
    \item 假设$\varphi_1,\varphi_2$满足命题并且$\varphi=(\varphi_1\to\varphi_2)$
    满足$\theta[\varphi]$, 此时也有$\theta[\varphi_1],\theta[\varphi_2]$, 从而
    \vspace{3pt}
    \[
        \overline{\eta\circ\theta}(\varphi)=
        \left(\overline{\eta\circ\theta}(\varphi_1)\to
        \overline{\eta\circ\theta}(\varphi_2)\right)=
        \left(\overline{\eta}\,\overline{\theta}(\varphi_1)\to
        \overline{\eta}\,\overline{\theta}(\varphi_2)\right)=
        \overline{\eta}\left(\overline{\theta}(\varphi_1)\to
        \overline{\theta}(\varphi_2)\right)=
        \overline{\eta}\,\overline{\theta}(\varphi).
    \]

    \vspace{3pt}
    \item 对任意的变元$x$, 假设$\varphi_1$满足命题并且$\varphi=\exists x\varphi_1$
    满足$\theta[\varphi]$, 此时也有$\theta/x[\varphi_1]$.

    \vspace{3pt}\hspace{2.17em}
    首先证明$(\eta/x)\circ(\theta/x)\overset{\displaystyle\varphi_1}{=}
    (\eta\circ\theta)/x$, 即任取$y\in\mathrm{at}(\varphi_1)$,
    证明$(\eta/x)\circ(\theta/x)(y)=(\eta\circ\theta)/x(y)$.

    \vspace{3pt}\hspace{2.17em}
    如果$y=x$, 那么显然两边都等于$x$;
    
    \hspace{2.17em}
    如果$y\neq x$, 则有$y\in\mathrm{at}(\varphi)$,
    根据$\theta[\varphi]$可知$x\notin\mathrm{var}\,\theta(y)$. 于是容易注意到
    $\eta/x\overset{\displaystyle\theta(y)}{=}\eta$, 所以
    \vspace{3pt}
    \[
        (\eta/x)\circ(\theta/x)(y)=\widehat{\eta/x}(\theta(y))=
        \widehat{\eta}(\theta(y))=(\eta\circ\theta)(y)=(\eta\circ\theta)/x(y).
    \]

    \vspace{3pt}\hspace{2.17em}
    刚刚的命题得证. 然后据此就有
    \begin{equation*}
        \overline{\eta\circ\theta}(\varphi)=
        \exists x\overline{(\eta\circ\theta)/x}(\varphi_1)=
        \exists x\overline{(\eta/x)\circ(\theta/x)}(\varphi_1)=
        \exists x\,\overline{\eta/x}\,\overline{\theta/x}(\varphi_1)=
        \overline{\eta}\left(\exists x\,\overline{\theta/x}(\varphi_1)\right)=
        \overline{\eta}\,\overline{\theta}(\varphi).
    \tag*{\qedsymbol}\end{equation*}
\end{enumerate}

% 又到晚上了.
% 向上 Löwenheim-Skolem 定理需要的引理太多了, 把所有概念和引理讲清楚需要巨大的篇幅.

% 这些证明几乎不可能记得住, 也很难快速读得懂, 然而似乎也没有真正的难点.
% 所以我大概不会再来读它们, 就让它们静静地躺在这无人问津的笔记中吧.

% 2026.03.01

\end{document}
